% huawei-cup-modeling working draft template
% IMPORTANT: This is NOT an official Huawei Cup / CPMCM submission template.
% Use it only when no current editable template is available.
% Before submission, migrate or verify formatting against the current official standard document.

\documentclass[12pt,a4paper]{ctexart}

\usepackage[left=2.5cm,right=2.5cm,top=2.5cm,bottom=2.5cm]{geometry}
\usepackage{amsmath,amssymb}
\usepackage{graphicx}
\usepackage{booktabs}
\usepackage{array}
\usepackage{longtable}
\usepackage{caption}
\usepackage{float}
\usepackage{placeins} % use \FloatBarrier selectively if floats drift across logical boundaries
\usepackage{enumitem}
\usepackage[hidelinks]{hyperref}

\setlength{\parindent}{2em}
\setlength{\parskip}{0pt}
\linespread{1.35}

\numberwithin{equation}{section}
\renewcommand{\theequation}{\thesection-\arabic{equation}}
\renewcommand{\thefigure}{\thesection-\arabic{figure}}
\renewcommand{\thetable}{\thesection-\arabic{table}}
\makeatletter
\@addtoreset{figure}{section}
\@addtoreset{table}{section}
\makeatother

\begin{document}

% WORKING DRAFT ONLY. Do not treat this layout as the current official competition template.

\begin{center}
    {\LARGE \textbf{论文题目}}\\[1em]
\end{center}

\begin{abstract}
在此撰写摘要。只填写已经由正文和结果台账支撑的方法、结果与结论，不补造数据。

\noindent\textbf{关键词：}关键词1；关键词2；关键词3
\end{abstract}

\section{问题重述}

\section{问题分析}

\section{模型假设}

\section{符号说明}

\section{模型建立与求解}

\section{结果分析与模型检验}

\section{模型评价}

\section{结论}

\begin{thebibliography}{99}
% \bibitem{ref1} ...
\end{thebibliography}

\appendix
\section{支撑材料说明}

\end{document}
